\newpage
\section*{Solucionario completo}

\noindent Cada respuesta incluye la clave y una justificación breve. En los ejercicios numéricos se indica la ecuación o el cálculo principal.

\subsection*{Bloque A. Fundamentos teóricos y conceptuales}

\sol{1}{C}{La demanda de trabajo no se deriva directamente del salario, sino de la demanda del bien o servicio que produce la empresa y de la productividad del trabajo. Por eso es una demanda derivada.}

\sol{2}{A}{La empresa competitiva contrata hasta que el ingreso adicional que genera el último trabajador, $MRP$, iguala el costo adicional de contratarlo, que es el salario.}

\sol{3}{D}{El ingreso marginal del producto del trabajo es $MRP=PMg\cdot IMg$. Si el mercado del producto es competitivo, $IMg=P$ y $MRP=VMP=P\cdot PMg$.}

\sol{4}{A}{Al aumentar el empleo, el producto marginal del trabajo disminuye por la ley de rendimientos marginales decrecientes; por eso la demanda de trabajo tiene pendiente negativa.}

\sol{5}{C}{La demanda de trabajo es más elástica cuando el bien tiene demanda elástica, el trabajo tiene sustitutos cercanos y el costo laboral representa una proporción importante de los costos totales.}

\sol{6}{C}{Un salario mayor eleva el costo relativo del trabajo y induce a sustituir trabajo por capital (efecto sustitución) y a reducir la producción (efecto escala). Ambos reducen la cantidad demandada de trabajo.}

\sol{7}{A}{La oferta de trabajo individual surge de la elección entre trabajo, que genera ingreso, y ocio, sujeto a una restricción de tiempo y a la necesidad de consumo.}

\sol{8}{A}{El efecto sustitución de un salario mayor encarece el ocio y incentiva a trabajar más. El efecto ingreso, en cambio, puede incentivar más ocio.}

\sol{9}{B}{Cuando el efecto ingreso domina al efecto sustitución, un salario mayor puede reducir la cantidad ofrecida de trabajo y generar una curva de oferta con pendiente negativa.}

\sol{10}{B}{El salario de reserva es el nivel mínimo de salario que compensa el costo de oportunidad del trabajo y induce a la persona a participar en el mercado laboral.}

\sol{11}{C}{Aunque la oferta individual puede tener tramos negativos, la oferta de trabajo del mercado generalmente tiene pendiente positiva: salarios más altos atraen más trabajadores.}

\sol{12}{C}{La elasticidad de la oferta de trabajo mide el cambio porcentual de la cantidad ofrecida ante un cambio porcentual del salario.}

\sol{13}{D}{El monopsonio es una estructura de mercado con un único comprador de un factor productivo, en este caso trabajo.}

\sol{14}{B}{Como el monopsonista enfrenta una oferta de trabajo con pendiente positiva, para contratar un trabajador adicional debe elevar el salario de todos los trabajadores; por eso el costo marginal del factor supera al salario.}

\sol{15}{A}{El monopsonista maximiza beneficios cuando el ingreso marginal del producto del trabajo iguala el costo marginal del factor trabajo: $MRP=CMgF$.}

\sol{16}{A}{El monopsonista restringe el empleo para pagar un salario menor. En comparación con la competencia, contrata menos trabajo y paga un salario inferior.}

\sol{17}{D}{La explotación monopsonística existe cuando el salario pagado es menor que el ingreso marginal del producto del trabajo: $w<MRP$.}

\sol{18}{B}{En un monopsonio, un salario mínimo fijado entre el salario de monopsonio y el competitivo puede aumentar el empleo, porque obliga a la empresa a pagar más y amplía la oferta de trabajo.}

\sol{19}{C}{El monopolio bilateral en el mercado de trabajo combina un único comprador de trabajo, normalmente una empresa, con un único vendedor organizado, normalmente un sindicato.}

\sol{20}{D}{En el monopolio bilateral no hay una solución única determinada por la oferta y la demanda; el salario y el empleo dependen del poder de negociación y del marco institucional.}

\sol{21}{A}{Los sindicatos pueden priorizar salarios, empleo, condiciones laborales o seguridad, según las preferencias de sus afiliados y el contexto de la negociación.}

\sol{22}{A}{La curva de contrato contiene las combinaciones de salario y empleo que son eficientes para el sindicato y la empresa; fuera de ella, ambas partes podrían mejorar.}

\sol{23}{C}{El poder de negociación depende de la facilidad para sustituir trabajo, la regulación, la cohesión sindical, la situación financiera de la empresa y el costo de un conflicto.}

\sol{24}{B}{La huelga es un mecanismo de presión que interrumpe la producción o el servicio y genera costos para la empresa, lo que afecta el resultado de la negociación.}

\sol{25}{A}{Las reglas de Hicks-Marshall señalan que la demanda de trabajo es más elástica cuando existen buenos sustitutos, la demanda del producto es elástica y el costo laboral es una proporción significativa de los costos.}

\sol{26}{C}{El cambio tecnológico que ahorra trabajo puede reducir la demanda por efecto sustitución, pero también puede aumentarla si reduce costos y expande la producción.}

\sol{27}{B}{Si el trabajo y el capital son sustitutos, un aumento del salario puede inducir a la empresa a reemplazar trabajo por capital y aumentar la demanda de capital.}

\sol{28}{B}{Los diferenciales compensatorios son diferencias salariales que compensan condiciones de trabajo desagradables, riesgosas o ubicaciones menos atractivas.}

\sol{29}{D}{La inversión en capital humano puede elevar la productividad del trabajo y, con ello, la demanda de trabajo y los salarios.}

\sol{30}{B}{En un mercado competitivo, el equilibrio se alcanza cuando la cantidad de trabajo demandada iguala la cantidad ofrecida al salario vigente.}

\subsection*{Bloque B. Aplicaciones cuantitativas y optimización}

\sol{31}{A}{En competencia, $70-2L=10+L$. Entonces $60=3L$, de donde $L=20$.}

\sol{32}{D}{Sustituyendo $L=20$ en la oferta: $w=10+20=30$.}

\sol{33}{B}{La oferta es $w=10+L$, por lo que el costo total del trabajo es $wL=10L+L^2$ y el costo marginal del factor es $CMgF=10+2L$. Igualando $MRP=CMgF$: $70-2L=10+2L$, de donde $L=15$.}

\sol{34}{D}{El monopsonista paga el salario de la oferta para $L=15$: $w=10+15=25$.}

\sol{35}{B}{$CMgF=10+2(15)=40$.}

\sol{36}{A}{$MRP=70-2(15)=40$. En el óptimo del monopsonista, $MRP=CMgF=40$.}

\sol{37}{D}{La explotación es $MRP-w=40-25=15$.}

\sol{38}{B}{El empleo competitivo es 20 y el de monopsonio es 15; la diferencia es 5.}

\sol{39}{D}{El salario competitivo es 30 y el de monopsonio es 25; la diferencia es 5.}

\sol{40}{A}{Con un salario mínimo de 30, la oferta de trabajo es $L=30-10=20$ y el $MRP$ en $L=20$ es 30. La empresa contrata 20 trabajadores.}

\sol{41}{D}{El salario mínimo que induce el empleo competitivo es el salario competitivo, 30. Por encima de ese nivel, el empleo empezaría a caer.}

\sol{42}{C}{$CMgF=10+2L$. Para $L=20$: $10+40=50$.}

\sol{43}{B}{La renta económica de los trabajadores es el área entre la oferta y el salario: $0.5(30-10)(20)=200$.}

\sol{44}{C}{En monopsonio, el área es $0.5(25-10)(15)=112.5$.}

\sol{45}{B}{La pérdida de eficiencia es el triángulo entre $MRP$ y la oferta desde $L=15$ hasta $L=20$: $0.5(20-15)(40-30)=25$.}

\sol{46}{B}{El costo total del trabajo es $wL=(10+L)L=10L+L^2$.}

\sol{47}{D}{Derivando $10L+L^2$ respecto de $L$: $CMgF=10+2L$.}

\sol{48}{C}{$MRP=PMg\cdot P=10\cdot 5=50$.}

\sol{49}{C}{$MRP=8\cdot 10=80$.}

\sol{50}{D}{La explotación es $MRP-w=60-40=20$.}

\sol{51}{A}{Con elasticidad $-2$, un aumento salarial de 10\% reduce el empleo en $2\cdot 10\%=20\%$.}

\sol{52}{C}{Con elasticidad 0.5, un aumento salarial de 10\% eleva la oferta en $0.5\cdot 10\%=5\%$.}

\sol{53}{B}{Elasticidad arco: $\frac{10/105}{2/21}=1$. La oferta es de elasticidad unitaria en ese tramo.}

\sol{54}{C}{Con $L=100-2w$, en $w=20$ se tiene $L=60$ y $dL/dw=-2$. La elasticidad es $-2(20/60)\approx -0.67$.}

\sol{55}{A}{Igualando $100-2w=2w-20$: $120=4w$, de donde $w=30$.}

\sol{56}{B}{Sustituyendo $w=30$: $L=100-60=40$ o $L=60-20=40$.}

\sol{57}{C}{Con salario mínimo 35, la demanda es $100-70=30$ y la oferta es $70-20=50$. El empleo efectivo es el mínimo: 30.}

\sol{58}{D}{El desempleo es la diferencia entre oferta y demanda: $50-30=20$.}

\sol{59}{D}{Un salario mínimo de 20 está por debajo del salario de equilibrio de 30, por lo que no es vinculante. El mercado permanece en el equilibrio competitivo y el empleo efectivo es 40.}

\sol{60}{A}{En un monopsonio, un salario mínimo igual al competitivo elimina la restricción monopsonística y eleva el empleo hasta el nivel competitivo.}

\sol{61}{C}{Un salario sindical por encima del competitivo reduce la cantidad demandada de trabajo y genera desempleo o menor empleo.}

\sol{62}{C}{En un monopolio bilateral, el resultado depende del poder de negociación relativo de la empresa y del sindicato.}

\sol{63}{A}{El salario negociado se ubica dentro de un rango delimitado por el salario de reserva de los trabajadores y el salario máximo que la empresa está dispuesta a pagar.}

\sol{64}{B}{La curva de contrato se obtiene cuando las curvas de indiferencia del sindicato son tangentes a las curvas de isobeneficio de la empresa.}

\sol{65}{C}{La masa salarial es $wL=w(100-2w)=100w-2w^2$. Derivando e igualando a cero: $100-4w=0$, de donde $w=25$.}

\subsection*{Bloque C. Casos con implicancias regulatorias y laborales}

\sol{66}{A}{La empresa de agua demanda trabajo porque demanda servicios de agua; la cantidad de trabajadores depende de la producción esperada y de su productividad.}

\sol{67}{B}{Si aumenta la demanda de servicios de agua, la empresa necesita más producción y, con ello, más trabajo, especialmente si la tecnología no cambia.}

\sol{68}{D}{La digitalización puede reemplazar la lectura manual y reducir la demanda de trabajo para esa tarea específica.}

\sol{69}{C}{Si la empresa es el único empleador de técnicos en la localidad, tiene poder de monopsonio en ese mercado laboral.}

\sol{70}{D}{El monopsonista paga un salario inferior al ingreso marginal del producto del trabajo; esa diferencia es la explotación monopsonística.}

\sol{71}{C}{Un salario mínimo entre el salario de monopsonio y el competitivo puede aumentar el empleo, porque la empresa enfrenta una oferta de trabajo más elástica en ese tramo.}

\sol{72}{B}{Si el salario mínimo supera el salario competitivo, la cantidad demandada de trabajo cae y el empleo disminuye.}

\sol{73}{D}{En la negociación colectiva, el resultado depende del poder de negociación, la organización sindical, la situación financiera y el costo de un conflicto.}

\sol{74}{D}{En un monopolio bilateral, el salario no queda determinado de forma única; se ubica en un rango y depende de la negociación.}

\sol{75}{A}{Si el sindicato prioriza el empleo, aceptará un salario menor a cambio de mantener o ampliar la contratación.}

\sol{76}{A}{La huelga impone costos a la empresa y al servicio, y por ello afecta el poder de negociación de las partes.}

\sol{77}{B}{El arbitraje obligatorio resuelve el conflicto mediante la decisión de un tercero, generalmente un árbitro o tribunal laboral.}

\sol{78}{A}{La capacitación eleva la productividad del trabajo y puede aumentar la demanda de trabajo y los salarios.}

\sol{79}{C}{La inversión que eleva la productividad del trabajo aumenta el $MRP$ y, con ello, la demanda de trabajo.}

\sol{80}{B}{La tercerización puede reducir la demanda de trabajo propio, pero aumentar la demanda de trabajadores especializados en las empresas contratistas; cambia la composición del empleo.}

\sol{81}{A}{Si la regulación reduce los ingresos de la empresa, esta puede ajustar costos y reducir la demanda de trabajo o moderar los salarios.}

\sol{82}{B}{El diferencial compensatorio retribuye la lejanía, las condiciones difíciles o los riesgos asociados al puesto.}

\sol{83}{D}{La oferta de trabajo en zonas remotas suele ser relativamente inelástica, porque hay menos trabajadores dispuestos a trasladarse.}

\sol{84}{C}{Si la oferta es inelástica, un aumento de la demanda de trabajo se traduce principalmente en salarios más altos y poco en más empleo.}

\sol{85}{B}{La demanda de trabajo es más elástica cuando el costo laboral es una proporción importante de los costos y existen sustitutos disponibles.}

\sol{86}{A}{Si la empresa no puede sustituir trabajo, la demanda es más inelástica y el empleo cae menos ante un aumento del salario.}

\sol{87}{D}{La negociación puede alcanzar un contrato eficiente si las partes negocian conjuntamente salario y empleo, en lugar de fijar solo el salario.}

\sol{88}{D}{Cuando coexisten poder de monopolio sindical y poder de monopsonio empresarial, el resultado es indeterminado y depende de la negociación.}

\sol{89}{C}{Un mayor costo de despido reduce la flexibilidad de la empresa para ajustar el empleo ante cambios en la demanda o en los costos.}

\sol{90}{D}{La evaluación de una política salarial debe considerar efectos sobre el empleo, los costos, la calidad del servicio, la sostenibilidad financiera y la equidad.}
